\input{../tex-inputs/latex-leseansicht-vorspann}

\section[Arthur Schnitzler an Gustav Schwarzkopf, 23. {[}4. 1911?{]}]{L04448 Arthur Schnitzler an Gustav Schwarzkopf, 23. {[}4. 1911?{]}}
\nopagebreak\mylabel{L04448v}
\rehead{ }\normalsize\beginnumbering\briefempfaengerindex{Schwarzkopf, Gustav@\textsc{Schwarzkopf, Gustav}!zzzSchnitzler, Arthur@\emph{von Arthur Schnitzler}!1911-04-231@{23. {[}4. 1911?{]}}|(be}
\toendnotes[C]{\smallbreak\pagebreak[2]}
\correspDesc{Versand  durch Arthur Schnitzler am 23. [4. 1911?] in Menton
\newline{}Übermittlung  am 23. 4. 1911 in Menton
\newline{}Erhalt  durch Gustav Schwarzkopf im Zeitraum [24. 4. 1911
                  – 28. 4. 1911?] in Wien}\toendnotes[C]{\smallbreak}
\Standort{DLA, A:Schnitzler, HS.1985.1.1897.}
\physDesc{Bildpostkarte, 176 Zeichen
\newline{}Handschrift: Bleistift, deutsche Kurrent
\newline{}Versand: Stempel: »\nobreak{}\oindex{Menton@\textbf{Menton}|pwk}Menton\hspace*{1em}Alpes Maritimes, 23 {[}4{]} {[}11{]}, \textcolor{gray}{2}{[}\textsuperscript{e}{]} 15\nobreak{}«.  }\toendnotes[C]{\smallbreak}\pstart{}{\pb}\textsc{Herrn Gustav Schwarzkopf}\pend{}\pstart{}\textsc{Wien I}\oindex{I., Innere Stadt@\textbf{I., Innere Stadt}|pw}\pend{}\pstart{}\textsc{Tiefer Graben 23}\oindex{Wien@\textbf{Wien}!I., Innere Stadt@\textbf{I., Innere Stadt}!Tiefer Graben 23@\textbf{Tiefer Graben 23}|pw}.\pend{}{\bigskip}
\pstart
           \noindent{}{\pb}\textcolor{gray}{\textbf{MENTON\oindex{Menton@\textbf{Menton}|pw}. – \begin{otherlanguage}{french}La Plage par un coup de Mer\end{otherlanguage}.}}\pend
           \vspace{1em}
\pstart
           \noindent{}{\pb}Herzliche Grüße; und auf \label{K_L04448-1v}\edtext{Wiederſehen Anfang Mai}{\lemma{\textnormal{\emph{Wiedersehen Anfang Mai}}}\Cendnote{\textnormal{Schnitzler kehrte am
                     3. 5. 1911 nach Wien\oindex{Wien@\textbf{Wien}|pwk} zurück und besuchte Schwarzkopf\pwindex{Schwarzkopf, Gustav 7.\,11.\,1853 Wien – 13.\,11.\,1939 ebd.@\textsc{Schwarzkopf, Gustav} (7.\,11.\,1853 Wien – 13.\,11.\,1939 ebd.)|pwk} fünf Tage später, vgl. A. S.: \emph{Tagebuch}, 8. 5. 1911.}}}\label{K_L04448-1}. \label{K_L04448-2v}\edtext{Wir fahren weg}{\lemma{\textnormal{\emph{Wir fahren weg}}}\Cendnote{\textnormal{Die Postkarte ist nicht datiert. Arthur und Olga Schnitzler\pwindex{Schnitzler, Olga 17.\,1.\,1882 Wien – 13.\,1.\,1970 Lugano@\textsc{Schnitzler, Olga} (17.\,1.\,1882 Wien – 13.\,1.\,1970 Lugano)|pwk}
                  verbrachten die Woche vom 16. 4. 1911 bis 24. 4. 1911 in Menton\oindex{Menton@\textbf{Menton}|pwk}. Der
                  Hinweis auf die nahe Abreise in Kombination mit der auf dem Poststempel lesbaren
                  Zahl »23« erlaubt die Datierung auf den 23. 4. 1911.}}}\label{K_L04448-2}
               nach München\oindex{München@\textbf{München}|pw}, \label{K_L04448-3v}\edtext{\textsc{Partenkirchen\oindex{Partenkirchen@\textbf{Partenkirchen}|pw}}}{\lemma{\textnormal{\emph{Partenkirchen}}}\Cendnote{\textnormal{Liesl Steinrück\pwindex{Steinrück, Elisabeth 19.\,11.\,1885 – 7.\,4.\,1920 Partenkirchen@\textsc{Steinrück, Elisabeth} (19.\,11.\,1885 – 7.\,4.\,1920 Partenkirchen)|pwk}, die unter Lungentuberkulose litt, wurde im Sanatorium Wigger\oindex{Dr. Wigger’s Kurheim Sanatorium für Innen-, Nervenkranke und Erholungsbedürftige@\textbf{Dr. Wigger’s Kurheim Sanatorium für Innen-, Nervenkranke und Erholungsbedürftige}|pwk} in Partenkirchen\oindex{Partenkirchen@\textbf{Partenkirchen}|pwk} behandelt. Die Schnitzlers\pwindex{Schnitzler, Olga 17.\,1.\,1882 Wien – 13.\,1.\,1970 Lugano@\textsc{Schnitzler, Olga} (17.\,1.\,1882 Wien – 13.\,1.\,1970 Lugano)|pwkv} besuchten sie auf der Hin- und auf der Rückreise für mehrere Tage (11. 4. 1911 bis 13. 4. 1911 und 27. 4. 1911 bis 1. 5. 1911).}}}\label{K_L04448-3},
                woher  ganz gute Nachrichten ko{\geminationm}en.\pend
           \pstart Ihr \spacefill\mbox{A. S.}\pend{}\selectlanguage{ngerman}\endnumbering\briefempfaengerindex{Schwarzkopf, Gustav@\textsc{Schwarzkopf, Gustav}!zzzSchnitzler, Arthur@\emph{von Arthur Schnitzler}!1911-04-231@{23. {[}4. 1911?{]}}|)be}\mylabel{L04448h}
\begin{anhang}
\end{anhang}\newcommand{\dateiname}{L04448}\newcommand{\titel}{Arthur Schnitzler an Gustav Schwarzkopf, 23. [4. 1911?]}\newcommand{\editorInnen}{Herausgegeben von Jahnke, SelmaMüller, Martin Anton}\input{../tex-inputs/latex-leseansicht-abspann}
